\subsection{EQUIVALENT SPECIES OF STRUCTURES}

Let \(\Sigma\) and \(\Theta\) be two species of structures, in the same theory \(\mathfrak{C}\) , having the same principal base sets \(x_{1},\ldots ,x_{n}\) . Let \(s,t\) be the generic structures of the species \(\Sigma,\Theta\) respectively. Suppose that the following conditions are satisfied:

\begin{enumerate}
\def\labelenumi{(\arabic{enumi})}
\item
  We have a procedure of deduction \(\mathbf{P}\{x_1, \ldots, x_n, s\}\) of a structure of species \(\Theta\) on \(x_1, \ldots, x_n\) from a structure of species \(\Sigma\) on \(x_1, \ldots, x_n\) .
\item
  We have a procedure of deduction \(\mathbb{Q}\{x_1, \ldots, x_n, t\}\) of a structure of species \(\Sigma\) on \(x_1, \ldots, x_n\) from a structure of species \(\Theta\) on \(x_1, \ldots, x_n\) .
\item
  The relation \(\mathbf{Q}\{x_1, \ldots, x_n, \mathbf{P}\{x_1, \ldots, x_n, s\}\} = s\) is a theorem in \(\mathfrak{G}_{\Sigma}\) , and the relation \(\mathbf{P}\{x_1, \ldots, x_n, \mathbf{Q}\{x_1, \ldots, x_n, t\}\} = t\) is a theorem in \(\mathfrak{G}_{\Theta}\) .
\end{enumerate}

The species of structures \(\Sigma\) and \(\Theta\) are then said to be equivalent by means of the procedures of deduction P and Q. In this case, for every theorem B \(\{x_1, \ldots, x_n, s\}\) in the theory \(\mathfrak{G}_{\Sigma}\) , the relation B \(\{x_1, \ldots, x_n, Q\}\) is a theorem in \(\mathfrak{G}_{\Theta}\) ; and conversely, for every theorem C \(\{x_1, \ldots, x_n, t\}\) in the theory \(\mathfrak{G}_{\Theta}\) , the relation C \(\{x_1, \ldots, x_n, P\}\) is a theorem in \(\mathfrak{G}_{\Sigma}\) .

If U is a structure of species \(\Sigma\) , the structure deduced from U by the procedure P is said to be equivalent to U. Criterion CST6 implies the following :

CST7. Let \(\mathcal{S},\mathcal{S}'\) be two structures of species \(\Sigma\) on the principal base sets \((\mathbf{E}_1,\ldots ,\mathbf{E}_n),(\mathbf{E}_1',\ldots ,\mathbf{E}_n')\) , respectively. Let \(\mathcal{S}_0,\mathcal{S}_0'\) be structures of species \(\Theta\) which are equivalent respectively to \(\mathcal{S}\) and \(\mathcal{S}'\) . In order that \((g_1,\dots,g_n)\) should be an isomorphism with respect to the structures \(\mathcal{S}_0\) and \(\mathcal{S}_0'\) , it is necessary and sufficient that \((g_1,\dots,g_n)\) should be an isomorphism with respect to the structures \(\mathcal{S}\) and \(\mathcal{S}'\) .

In practice, we make no distinction between the theories \(\mathfrak{G}_{\Sigma}\) and \(\mathfrak{G}_{\Theta}\) of two equivalent species of structures.

\minorheading{Examples}

* (1) Let \(\Sigma\) be the species of commutative group structures; \(\Sigma\) has a single (principal) base set A, and its generic structure consists of a single letter F; the typical characterization of \(\Sigma\) is \(F \in \mathfrak{P}((A \times A) \times A)\) , and we denote the axiom of \(\Sigma\) by R \{A, F\}. This axiom implies in particular that F is the graph of a function (the ``law of composition'' of the group; cf.~no. 4, Example 2). In the theory \(\mathfrak{G}_{\Sigma}\) (where \(\mathfrak{G}\) denotes the theory of sets) we define a term M\{A, F\} which is a functional graph in \(\mathfrak{P}((Z \times A) \times A)\) and satisfies the following relation B\{M, A, F\}:

\((\forall x)(\forall y)(\forall n)((x \in A \text{ and } y \in A \text{ and } n \in \mathbf{Z})\)

\[
\Longrightarrow (\mathbf {M} (n, \mathrm{F} (x, y)) = \mathrm{F} (\mathbf {M} (n, x), \mathrm{M} (n, y))))
\]

and \((\forall x)(\forall m)(\forall n)((x\in A\) and \(m\in \mathbf{Z}\) and \(n\in \mathbf{Z})\)

\[
\Rightarrow (\mathrm{M} (m + n, x) = \mathrm{F} (\mathrm{M} (m, x), \mathrm{M} (n, x))))
\]

and \((\forall x)(\forall m)(\forall n)((x\in A\) and \(m\in \mathbf{Z}\) and \(n\in \mathbf{Z})\)

\[
\Longrightarrow (\mathbf {M} (m, \mathbf {M} (n, x)) = \mathbf {M} (m n, x)))
\]

and

\[
(\forall x) ((x \in A) \Longrightarrow (M (1, x) = x)).
\]

(``multiplication of an element of A by an integer'').

Consider the species \(\Theta\) of \(\mathbf{Z}\) -module structures, which has a single principal base set A, with \(\mathbf{Z}\) as auxiliary set, and whose generic structure contains two letters G, L, with the typical characterization

\[
\mathbf {G} \in \mathfrak {P} ((\mathbf {A} \times \mathbf {A}) \times \mathbf {A}) \quad \text { and } \quad \mathbf {L} \in \mathfrak {P} ((\mathbf {Z} \times \mathbf {A}) \times \mathbf {A})
\]

and the axiom

\[
\text{``}R\{A,G\}\text{ and }(L\text{ is a functional graph})\text{ and }B\{L,A,G\}\text{''.}
\]

It is immediately verified that the terms F, M constitute a procedure of deduction of a structure of species \(\Theta\) from a structure of species \(\Sigma\) , and that the term G is a procedure of deduction of a structure of species \(\Sigma\) from a structure of species \(\Theta\) . Furthermore, the condition (3) above is trivially satisfied. We may therefore say that the species of commutative group structures and the species of Z-module structures are equivalent.

\begin{enumerate}
\def\labelenumi{(\arabic{enumi})}
\setcounter{enumi}{1}
\tightlist
\item
  Let \(\Sigma\) be the species of topological structures (no. 4, Example 3), let A be the (principal) base set, and let V be the generic structure of \(\Sigma\) . Consider the relation
\end{enumerate}

\[
x \in A \text {   and   } X \subset A \text {   and   } (\forall U) ((U \in V \text {   and   } x \in U) \Longrightarrow (X \cap U \neq \emptyset)).
\]

This relation has a graph \(P \subset \mathfrak{P}(A) \times A\) with respect to the pair \((X, x)\) ; \(P\{A, V\}\) is a term called ``the set of all pairs \((X, x)\) such that x lies in the closure of X with respect to the topology V''. We can then prove (cf.~General Topology, Chapter I, §1) that the following relations are theorems in \(G_{\Sigma}\) :

\[
\begin{array}{c} \mathrm{P} (\phi) = \phi , \\ (\forall \mathrm{Y}) ((\mathrm{Y} \subset \mathrm{A}) \Longrightarrow (\mathrm{Y} \subset \mathrm{P} (\mathrm{Y}))), \\ (\forall \mathrm{Y}) ((\mathrm{Y} \subset \mathrm{A}) \Longrightarrow (\mathrm{P} (\mathrm{P} (\mathrm{Y})) = \mathrm{P} (\mathrm{Y}))), \\ (\forall \mathrm{Y}) (\forall \mathrm{Z}) ((\mathrm{Y} \subset \mathrm{A} \text {and} \mathrm{Z} \subset \mathrm{A}) \Longrightarrow (\mathrm{P} (\mathrm{Y} \cup \mathrm{Z}) = \mathrm{P} (\mathrm{Y}) \cup \mathrm{P} (\mathrm{Z}))). \end{array}
\]

Consider the species of structures \(\Theta\) , having a single (principal) base set A, whose generic structure consists of a single letter W, which has as typical characterization \(W \in \mathfrak{P}(\mathfrak{P}(A) \times A)\) and as axiom

\[
\begin{array}{c} \mathrm{W} (\emptyset) = \emptyset \quad \text {and} \quad (\forall \mathrm{Y}) ((\mathrm{Y} \subset \mathrm{A}) \Longrightarrow (\mathrm{Y} \subset \mathrm{W} (\mathrm{Y}))) \\ \text {and} \quad (\forall \mathrm{Y}) ((\mathrm{Y} \subset \mathrm{A}) \Longrightarrow (\mathrm{W} (\mathrm{W} (\mathrm{Y})) = \mathrm{W} (\mathrm{Y}))) \\ \text {and} \quad (\forall \mathrm{Y}) (\forall \mathrm{Z}) ((\mathrm{Y} \subset \mathrm{A} \quad \text {and} \quad \mathrm{Z} \subset \mathrm{A}) \Longrightarrow (\mathrm{W} (\mathrm{Y} \cup \mathrm{Z}) = \mathrm{W} (\mathrm{Y}) \cup \mathrm{W} (\mathrm{Z}))). \end{array}
\]

Consider also the relation

\[
\mathrm{U} \subset \mathrm{A} \quad \text { and } \quad (\forall x) ((x \in \mathrm{U}) \Longrightarrow x \notin \mathrm{W} (\mathrm{A} - \mathrm{U})).
\]

The set of all \(U \in \mathfrak{P}(A)\) which satisfy this relation is a subset \(Q\{A, W\}\) of \(\mathfrak{P}(A)\) . We can then prove (General Topology, Chapter I, §1, Exercise 10) that the following relations are theorems in \(G_{\Theta}\) :

\[
\begin{array}{c} \mathrm{A} \in \mathrm{Q}, \\ (\forall \mathrm{M}) ((\mathrm{M} \subset \mathrm{Q}) \Longrightarrow \left(\left(\bigcup_ {\mathrm{X} \in \mathrm{M}} \mathrm{X}\right) \in \mathrm{Q}\right), \end{array}
\]

\[
(\forall \mathrm{X}) (\forall \mathrm{Y}) ((\mathrm{X} \in \mathrm{Q} \text {   and   } \mathrm{Y} \in \mathrm{Q}) \Longrightarrow ((\mathrm{X} \cap \mathrm{Y}) \in \mathrm{Q})).
\]

Thus the terms \(P\{A,V\}\) and \(Q\{A,W\}\) satisfy conditions (1) and (2) above, and it is easily seen that they also satisfy condition (3). The species of structures \(\Sigma\) and \(\Theta\) are therefore equivalent, and we therefore consider every structure of species \(\Theta\) as a topology, namely that which corresponds to it under the procedure of deduction \(Q\{A,W\}\). *

